\documentclass[perfil=artigo]{undftex}
\usepackage{longtable}
\addbibresource{tests/apud.bib}
\ExplSyntaxOn
\cs_new:Npn \linhasteste { \prg_replicate:nn {65} { Dado~de~teste & Conteúdo~de~linha~para~teste~multipágina. \\ } }
\cs_new:Npn \citacaoteste { \prg_replicate:nn {24} { Este~é~um~texto~fictício~para~testar~a~quebra~de~uma~citação~extensa.~Não~representa~transcrição~de~obra~real.~\par } }
\ExplSyntaxOff
\begin{document}
\textual\singlespacing
\noindent Teste: um bloco curto deve caber no espaço disponível.
\vspace*{.58\textheight}
\begin{undfpseudocodigo}{TESTE(V)}{curto}{Título descritivo suficientemente extenso para verificar a quebra automática de linha sem ultrapassar a largura do texto}
Entrada: vetor V
Inicio
  Retorne V
Fim
\end{undfpseudocodigo}
\ifnum\value{page}=1\else\errmessage{Bloco curto foi antecipado para outra pagina}\fi
\clearpage
\noindent Teste: uma citação curta deve caber no espaço disponível.
\vspace*{.69\textheight}
\begin{citacaolonga}[9]{consultada}
Este trecho fictício verifica a composição da citação em bloco. Não é uma
transcrição de obra real. Seu tamanho foi escolhido para caber no espaço
restante da página, sem a reserva de uma fração fixa da altura textual.
\end{citacaolonga}
\ifnum\value{page}=2\else\errmessage{Citacao curta foi antecipada para outra pagina}\fi
\clearpage
\section{Pseudocódigo multipágina}
\label{obj:pseudocodigo-inicio}
\begin{undfpseudocodigolongo}{PERCORRER(V)}{longo}{Percurso extenso usado somente para verificar a continuação do quadro entre páginas}
Entrada: vetor V
Inicio
  Execute etapa 01
  Execute etapa 02
  Execute etapa 03
  Execute etapa 04
  Execute etapa 05
  Execute etapa 06
  Execute etapa 07
  Execute etapa 08
  Execute etapa 09
  Execute etapa 10
  Execute etapa 11
  Execute etapa 12
  Execute etapa 13
  Execute etapa 14
  Execute etapa 15
  Execute etapa 16
  Execute etapa 17
  Execute etapa 18
  Execute etapa 19
  Execute etapa 20
  Execute etapa 21
  Execute etapa 22
  Execute etapa 23
  Execute etapa 24
  Execute etapa 25
  Execute etapa 26
  Execute etapa 27
  Execute etapa 28
  Execute etapa 29
  Execute etapa 30
  Execute etapa 31
  Execute etapa 32
  Execute etapa 33
  Execute etapa 34
  Execute etapa 35
  Execute etapa 36
  Execute etapa 37
  Execute etapa 38
  Execute etapa 39
  Execute etapa 40
  Execute etapa 41
  Execute etapa 42
  Execute etapa 43
  Execute etapa 44
  Execute etapa 45
  Execute etapa 46
  Execute etapa 47
  Execute etapa 48
  Execute etapa 49
  Execute etapa 50
  Execute etapa 51
  Execute etapa 52
  Execute etapa 53
  Execute etapa 54
  Execute etapa 55
  Execute etapa 56
  Execute etapa 57
  Execute etapa 58
  Execute etapa 59
  Execute etapa 60
  Execute etapa 61
  Execute etapa 62
  Execute etapa 63
  Execute etapa 64
  Execute etapa 65
  Retorne V
Fim
\end{undfpseudocodigolongo}
\label{obj:pseudocodigo-fim}
\clearpage
\section{Tabela multipágina}
\label{obj:tabela-inicio}
\begin{longtable}{p{.25\textwidth}p{.60\textwidth}}
\caption{Tabela de teste multipágina}\label{tab:longa}\\
\textbf{Item} & \textbf{Descrição}\\\hline\endfirsthead
\textbf{Item} & \textbf{Descrição (continuação)}\\\hline\endhead
\linhasteste
\end{longtable}
\label{obj:tabela-fim}
\undffonte{dados fictícios para teste}
\section{Figura e citação extensa}
\undffigura[.4\textwidth]{identidade/logo_horizontal.png}{Figura de teste com proporções preservadas}{objeto}
\label{obj:citacao-inicio}
\begin{citacaolonga}{consultada}
\citacaoteste
\end{citacaolonga}
\label{obj:citacao-fim}
\printundfreferences
\end{document}
